% TCC — Resultados e Discussão
\section{Resultados e Discussão}\label{sec:resultados}

Esta seção registra inicialmente os resultados do piloto usados para auditar e
revisar o mecanismo de priorização. Os números ainda não constituem a avaliação
externa final do FIS revisado. Resultados negativos e casos que contradigam as
hipóteses são preservados porque explicam as decisões posteriores do projeto.

\subsection{Caracterização do Piloto}

Foram executadas campanhas de nível 2 contra VAmPI, Spring Petclinic REST,
crAPI e RealWorld. No conjunto canônico, 958 requisições planejadas foram
executadas, das quais 284 continham combinações de duas mutações. A
deduplicação estrita produziu 265 assinaturas: 7 na VAmPI, 47 na Petclinic, 198
na crAPI e 13 na RealWorld. Para a avaliação de prioridade, essas assinaturas
foram reunidas em 161 famílias de comportamento.

O nível 2 ampliou a variedade das manifestações observadas, especialmente na
crAPI, onde 37 assinaturas novas envolveram combinações L2. No entanto, o
piloto não produziu uma nova família de anomalia de prioridade alta exclusiva
desse nível. O resultado não torna o L2 irrelevante: divergências de contrato,
respostas combinadas e variações de robustez também fazem parte do objeto de
estudo, mesmo quando não indicam vulnerabilidade grave.

\subsection{Construção da Referência Cega}

Três julgadores agentes avaliaram as 161 famílias sem acesso ao rótulo do
analisador ou aos resultados do FIS. Houve concordância exata sobre validade em
144 famílias (89,4\%). As mesmas 144 foram consideradas suportadas pelos três
julgadores, com concordância exata de prioridade em 136 casos (94,4\%). Sete
casos sem maioria de prioridade foram submetidos a um quarto julgador. A
referência final manteve todas as famílias como suportadas e distribuiu-as em
99 LOW, 59 MEDIUM e 3 HIGH.

As três famílias HIGH corresponderam a duas falhas 500 repetidas em baselines
válidos da Petclinic e a uma operação autenticada da crAPI que aceitou
credencial ausente e retornou conteúdo funcional. A referência é descrita como
operacional e assistida por agentes, não como verdade humana absoluta.

\subsection{Avaliação da Priorização Fuzzy}

A auditoria do FIS v2 mostrou que o status HTTP dominava a decisão. A ablação
do tempo e a ablação do sinal de tamanho não alteraram a faixa de nenhuma das
161 famílias. O v2 e o baseline baseado apenas em status obtiveram a mesma
acurácia, 0,298, o mesmo macro-F1, 0,254, e o mesmo erro ordinal médio, 0,702.
O v2 apresentou correlação de Spearman de 0,710 com a referência.

A matriz de confusão explicou a baixa acurácia: 67 famílias LOW foram elevadas
a MEDIUM e 43 famílias MEDIUM foram elevadas a HIGH. O modelo
superpriorizava rejeições seguras, status 4xx não documentados e falhas 5xx
localizadas. Em particular, uma resposta 4xx não documentada pode representar
contract drift e ainda assim constituir rejeição segura da entrada exercitada;
as duas propriedades precisam ser preservadas sem receber automaticamente a
mesma prioridade de uma operação válida indisponível.

Esse resultado motivou o FIS v3. Seus antecedentes são contradição
comportamental, força da evidência e exposição de controle de acesso. Tempo e
tamanho deixaram de integrar o núcleo da decisão. Repetição após deduplicação
fortalece a evidência de falhas de servidor; credencial ausente e inválida são
tratadas separadamente; e a conformidade da resposta é recalculada contra o
contrato, sem usar o rótulo final do analisador como antecedente.

Aplicado às mesmas 161 famílias de desenvolvimento, o código do v3 reproduziu
todas as faixas da referência, com macro-F1 1,000 e erro ordinal zero. Esse
resultado comprova que a implementação expressa a rubrica formulada, mas não
estima generalização: as regras foram definidas depois da análise desta
amostra. Uma quinta API ou um conjunto futuro congelado antes da rotulagem é
necessário para a avaliação confirmatória.

\subsection{Casos Representativos e Falsos Positivos}

Os desacordos do v2 mostraram três padrões importantes. Primeiro, rejeições
404 de identificadores gerados não comprovam falha da operação e tendem a LOW.
Segundo, um baseline POST cujo corpo satisfaz o schema OpenAPI, mas recebe 400
por restrições não declaradas pelo contrato, evidencia contract drift de
entrada e tende a MEDIUM. Terceiro, um erro 500 isolado possui evidência mais
fraca que o mesmo comportamento repetido em diferentes casos; por isso, a
recorrência atua como modificador de prioridade no v3.

\subsection{Discussão Geral}

O principal resultado do piloto não é que o FIS v2 tenha superado um baseline,
mas que a instrumentação permitiu demonstrar por que ele não o superava. A
revisão substituiu sinais fáceis de medir, porém pouco discriminativos, por
fatos mais próximos do construto. A contribuição deve ser apresentada como um
processo auditável de priorização de anomalias técnicas, incluindo contract
drift e robustez, e não como estimador de impacto real ou gravidade de
segurança.

Permanecem como ameaças a quantidade de APIs, a dependência do OpenAPI como
oráculo parcial, a correlação entre julgadores agentes e o ajuste do v3 no
mesmo conjunto usado para diagnosticar o v2. As campanhas definitivas devem
congelar código, regras e protocolo antes da coleta confirmatória.
